\chapter*{Interlude: The Life Story of the Number Seven}
\addcontentsline{toc}{chapter}{Interlude: The Life Story of the Number Seven}
\markboth{INTERLUDE}{INTERLUDE}

\begin{oberflaeche}
Let us accompany a single number through all of its homes.

The seven is born into the world with twenty-four places -- earlier
is impossible; the world before it ends at five. There it is called
$013$. It is a cramped apartment: three digits, all of them
carrying something, nothing is decoration.

Then it moves, world after world, always value-preserving -- it
wants to remain the same seven, after all. And with every move it
has to furnish itself anew: in the world with one hundred and
twenty places it is called $0012$, in the next one $00011$, then
$000010$. Lay the moves side by side and you see a restless
pattern -- the digits wander, swap, shrink. No two moves are alike,
and nothing in the spelling betrays at first glance that it is
always the same number.

Until the world with forty thousand three hundred and twenty
places. There something new happens: for the first time the seven
fits entirely into the last digit -- it is now called $0000007$,
and you can see its value again. From this move on, nothing ever
changes again. Each further world merely hangs one more silent zero
in front: $00000007$, $000000007$, and so on forever. The seven has
arrived; we say: it has reached its eigenhorizon.

Every number has this biography: a restless youth, in which each
larger space means a fresh negotiation of identity, and a quiet
adulthood, from which point on more space is simply more space.
Large numbers have a long youth -- the number seven hundred and
twenty, for instance, keeps being rebuilt until world number seven
hundred and twenty, and along that long road it temporarily wears
shapes nobody would credit it with: in the ninth world it is simply
a single one in the fourth position from the right, because ten
times nine times eight is seven hundred and twenty. Exactly this
double nature -- unwieldy by value, plain by shape -- later turned
out to be the key to one of the main theorems.

Only one thing needs to be remembered from all this: in the Hori
space, ``the same number'' is not a matter of course but a
decision. You can hold the value fixed and sacrifice the shape, or
the other way round. How much that decision costs is no longer a
matter of taste but a theorem -- but that belongs in the next
chapter.
\end{oberflaeche}
